\section*{第 \ref{ch:g9:kinetic-energy-safety} 章 --- 动能与道路安全}

\begin{solution}{exo:g9:kinetic-energy-safety:1}
焦耳（\unit{J}）：几焦耳就能把这本书举到
书架上；心脏每跳一次大约耗费一焦耳。$E_k = \frac12 m
v^2$ --- $m$ 以千克为单位，$v$ 以米每秒为单位，$E_k$ 以
焦耳为单位。
\end{solution}

\begin{solution}{exo:g9:kinetic-energy-safety:2}
野兔：$0.5 \times 2 \times 100 = \qty{100}{J}$。运动员：$0.5
\times 80 \times 25 = \qty{1000}{J}$。
\end{solution}

\begin{solution}{exo:g9:kinetic-energy-safety:3}
质量增加一半：倍数为 $1.5$。速度增加一半：倍数为 $1.5^2
= 2.25$。速度这一项占优——它总是如此，因为它带着
平方。
\end{solution}

\begin{solution}{exo:g9:kinetic-energy-safety:4}
反应距离——一秒内行驶的路程，
与 $v$ 成正比增长；制动距离——
即 $E_k$ 的“退役”，随 $v^2$ 增长。
\end{solution}

\begin{solution}{exo:g9:kinetic-energy-safety:5}
在 \qty{50}{km/h} 时：$13$ 占 \qty{27}{m} 的一半左右。在
\qty{130}{km/h} 时：$85$ 占 \qty{121}{m} 的约百分之七十。
平方项的增长快于正比项：
高速时，制动几乎主导了整个停车距离。
\end{solution}

\begin{solution}{exo:g9:kinetic-energy-safety:6}
转化为制动盘上的热量——长下坡之后的琥珀色光晕，
以及制动过热的焦味。能量只是被传递，
从不会被抹去：这正是早先“能量”一章中的链式命题，
依然成立。
\end{solution}

\begin{solution}{exo:g9:kinetic-energy-safety:7}
$36 \div 3.6 = \qty{10}{m}$；接着是 \qty{20}{m}；再
\qty{30}{m}——与速度成正比推进，
正如一秒的路程理应如此。
\end{solution}

\begin{solution}{exo:g9:kinetic-energy-safety:8}
你的 $E_k$ 由你的质量和汽车的速度决定——没有
任何安全带能改变它。安全带使\emph{退役}变得文明些：
让身体在更长的距离和时间内温和地停下，
而不是一下子撞向仪表盘。
\end{solution}

\begin{solution}{exo:g9:kinetic-energy-safety:9}
制动距离与速度的平方成正比：在 \qty{60}{km/h}
（为两倍），$8 \times 4 = \qty{32}{m}$；在 \qty{90}{km/h}
（为三倍），$8 \times 9 = \qty{72}{m}$。
\end{solution}

\begin{solution}{exo:g9:kinetic-energy-safety:10}
在 \qty{30}{km/h} 时：反应距离 $\approx \qty{8.3}{m}$，制动
距离 $13 \times (30/50)^2 \approx \qty{4.7}{m}$——停车距离共约
\qty{13}{m}，对比 \qty{50}{km/h} 时的 \qty{27}{m}：不到一半的路程。
能量方面：$(30/50)^2 = 0.36$——仅约碰撞能量的三分之一。
两者都有说服力；海报上最合适的标语，通常是道路那一句——
“以 30 行驶，你会在孩子前面停下；
以 50 行驶，你到达孩子身边时仍在移动”。
\end{solution}

\begin{solution}{exo:g9:kinetic-energy-safety:11}
只将制动部分加倍——雨天拉长的是
制动的“退役”，而不是驾驶员的反应：在 \qty{50}{km/h} 时，$14 + 26 =
\qty{40}{m}$；在 \qty{90}{km/h} 时，$25 + 80 = \qty{105}{m}$。
反应的那一秒，干路面与湿路面并无二致。
\end{solution}

\begin{solution}{exo:g9:kinetic-energy-safety:12}
卡车：$v = 25$ \unit{m/s}；$E_k = 0.5 \times 20000 \times
625 = \qty{6.25e6}{J}$。能量相当的小汽车：$0.5 \times 1000 \times
v^2 = \num{6.25e6}$，得 $v^2 = 12500$，$v \approx
\qty{112}{m/s} \approx \qty{400}{km/h}$——没有任何道路车辆
能接近这一速度。卡车以高速公路速度行驶时，携带的是赛车级的
能量，却只配备货车的制动器：一旦制动器在长下坡上过热失灵，
只有碎石避险车道能安全地让兆焦级的能量“退役”。
小汽车从不需要这类车道；它们的能量在普通制动器中
便已耗散。
\end{solution}

\begin{solution}{pb:g9:kinetic-energy-safety:1}
\textbf{1.} $50 \div 3.6 = \qty{13.9}{m/s}$；$60 \div 3.6 =
\qty{16.7}{m/s}$。
\textbf{2.} \qty{13.9}{m} 与 \qty{16.7}{m}。
\textbf{3.} \qty{13}{m}；在时速 60 时，$13 \times (60/50)^2 =
13 \times 1.44 \approx \qty{18.7}{m}$。
\textbf{4.} 停车距离：$26.9$，对比 \qty{35.4}{m}——海报上的
空白处：约多出 \emph{八米半}，
即比时速 50 司机停下的地方还远两个车身。
\textbf{5.} 时速 50 时：约 \qty{97}{kJ}。时速 90 时（$v =
\qty{25}{m/s}$）：$0.5 \times 1000 \times 625 =
\qty{312}{kJ}$。
\textbf{6.} $97 \div 9.8 \approx \qty{10}{m}$；$312 \div 9.8
\approx \qty{32}{m}$。
\textbf{7.} 十米相当于三层楼房；三十二米相当于
十层楼。草稿写道：“时速 50 的撞击，等同于从三楼
坠落；时速 90，则是从十楼坠落。请选好你的窗口。”
\textbf{8.} 这些高度如实反映了必须被“退役”的能量——
这个数字不会因任何工程改进而改变。
吸能区与安全带改变的是“退役”的
\emph{方式}：以米计的温和制停，而非
以厘米计的猛烈制停——这正是坠落类比
夸大现代汽车中伤害的原因，但它所描绘的能量
依然真实。
\textbf{9.} 每秒 \qty{25}{m} 的盲区；两秒的
一瞥，\qty{50}{m}——还没等到反应那一秒开始，
就已驶过半个球场的距离。
\textbf{10.} $65 + 25 = \qty{90}{m}$——疲惫的司机
比清醒司机的总距离多出一个半公交车长。
\textbf{11.} 在\emph{低速}时上述说法偏差最大：此时
平方项的制动份额很小，而反应路段
占了停车距离的大部分——在城市里，那一秒钟
\emph{正是}危险所在。（高速时，制动占主导——
但每秒 \qty{25}{m} 在那里也绝非小事。）
\textbf{12.} 例如：“速度多 10，停车距离多
10 米——速度的影响出现两次，其中一次是平方。
撞击就是坠落：能量与速度的平方成正比增长。
一秒钟的一瞥，就是 \qty{25}{m} 的盲区——
注意力才是最短的制动距离。”
\end{solution}
